\section{TỔNG QUAN ĐỀ TÀI}
\label{sec:chuong1}

\subsection{Bối cảnh và nhu cầu tối ưu quy trình lấy hàng trong kho}
Lấy hàng là quá trình thu gom sản phẩm tại các vị trí lưu trữ để đáp ứng yêu cầu của đơn hàng. Nhân viên di chuyển từ điểm tập kết đến các vị trí cần lấy, thu gom sản phẩm và quay về điểm tập kết khi kết thúc chuyến. Trong báo cáo, điểm tập kết được gọi là depot; một mã hàng phân biệt sản phẩm được gọi là Stock Keeping Unit (SKU). Một đơn có thể gồm nhiều mã hàng nằm ở các vị trí khác nhau, còn một chuyến có thể phục vụ nhiều đơn nếu tổng tải không vượt sức chứa xe đẩy.

Khi nhiều đơn cùng chờ xử lý, cách tổ chức lấy hàng ảnh hưởng đến cả quãng đường di chuyển và thời điểm hoàn thành. Phục vụ từng đơn riêng lẻ có thể làm nhân viên đi lặp lại các đoạn đường. Gom đơn cho phép chia sẻ lộ trình, nhưng cũng có thể kéo dài chuyến và khiến đơn có hạn gần phải chờ lâu hơn. Khi có nhiều nhân viên, việc phân công và sắp xếp chuyến còn quyết định mức cân bằng công việc giữa các nhân viên.

Tổng quan của de Koster và cộng sự \cite{dekoster2007} hệ thống hóa các quyết định về thiết kế và điều hành hoạt động lấy hàng, trong đó có định tuyến và gom đơn. Với mô hình được xét trong đề tài, yêu cầu đặt ra là xây dựng kế hoạch phục vụ đủ đơn, tuân thủ sức chứa và các lối đi của kho, đồng thời cân đối quãng đường, thời gian hoàn tất và trễ hạn. Đây là cơ sở để lựa chọn một bài toán tối ưu có sự kết hợp giữa nhiều quyết định vận hành.

\subsection{Bài toán tích hợp gom đơn, định tuyến và lập lịch}
Đề tài xét bài toán tích hợp ba quyết định: gom đơn xác định các đơn được phục vụ trong cùng một chuyến; định tuyến xác định thứ tự ghé các vị trí lấy hàng và đường đi giữa chúng; lập lịch xác định nhân viên phụ trách cùng thứ tự thực hiện các chuyến của mỗi người. Một nghiệm hoàn chỉnh phải thể hiện được cả ba quyết định này.

Các quyết định có quan hệ chặt chẽ. Thêm một đơn vào chuyến làm thay đổi tập vị trí cần ghé và có thể làm tăng thời lượng chuyến. Các chuyến tiếp theo của cùng nhân viên khi đó phải bắt đầu muộn hơn, ảnh hưởng đến thời điểm hoàn thành những đơn khác. Ngược lại, thay đổi nhân viên phụ trách hoặc thứ tự chuyến có thể giảm trễ mà không cần thay đổi cách gom đơn.

Chẳng hạn, hai đơn có sản phẩm gần nhau nhưng hạn hoàn thành khác nhau có thể được gom để chia sẻ lộ trình. Tuy nhiên, trong mô hình của đề tài, đơn chỉ hoàn thành khi chuyến chứa nó quay về depot; vì vậy đơn gấp có thể phải chờ việc lấy hàng của đơn còn lại. Tách đơn gấp thành chuyến riêng có thể giảm trễ nhưng làm tăng quãng đường. Ví dụ này cho thấy một thay đổi có lợi cho một chỉ tiêu chưa chắc có lợi cho các chỉ tiêu còn lại.

Đề tài sử dụng hạn hoàn thành mềm: một đơn được phép hoàn thành sau hạn, nhưng phần thời gian vượt hạn được tính vào chi phí. Chất lượng nghiệm được đánh giá bằng tổng có trọng số của ba thành phần đã chuẩn hóa: tổng quãng đường, thời gian hoàn tất toàn bộ công việc (makespan) và tổng độ trễ. Số đơn trễ được theo dõi như một chỉ tiêu bổ sung. Khi phân tích thực nghiệm, cần xem cả mục tiêu tổng hợp lẫn từng thành phần để nhận diện các đánh đổi.

\subsection{Các hướng giải quyết liên quan}
Các nghiên cứu về lấy hàng có thể tập trung vào một quyết định riêng hoặc xét nhiều quyết định đồng thời. Gademann và van de Velde \cite{gademann2005} nghiên cứu gom đơn nhằm giảm tổng thời gian di chuyển trong kho có các lối song song. Ratliff và Rosenthal \cite{ratliff1983} xây dựng phương pháp định tuyến cho một trường hợp kho chữ nhật với lối nối ở hai đầu dãy. Hai công trình cung cấp cơ sở tham khảo cho việc tổ chức nhóm đơn và khai thác cấu trúc đường đi.

Ở hướng tích hợp, Briant và cộng sự \cite{briant2023} xét bài toán gom đơn, định tuyến và sắp thứ tự có hạn hoàn thành, sử dụng mô hình quy hoạch tuyến tính và phương pháp sinh cột để xây dựng cận dưới và cận trên. Đề tài tham khảo cách tiếp cận tích hợp, nhưng sử dụng mô hình có hạn mềm và mục tiêu gồm ba thành phần. Vì khác biệt về mô hình và tiêu chí đánh giá, kết quả được đối chiếu giữa các phương pháp cài đặt trong đề tài trên cùng dữ liệu và cùng bộ đánh giá.

Về tìm kiếm nghiệm, Shaw \cite{shaw1998} trình bày Large Neighborhood Search (LNS), thay đổi một phần nghiệm bằng cách loại bỏ và chèn lại các phần tử. Ropke và Pisinger \cite{ropke2006} phát triển Adaptive Large Neighborhood Search (ALNS), trong đó tần suất sử dụng các toán tử phụ thuộc vào hiệu quả đã quan sát. Mladenović và Hansen \cite{mladenovic1997} đề xuất Variable Neighborhood Search (VNS), dựa trên việc thay đổi có hệ thống cấu trúc lân cận. Các khung này là cơ sở để xây dựng những chiến lược cải thiện nghiệm khác nhau cho bài toán kho hàng.

\begin{table}[H]
\centering\small
\begin{tabularx}{\linewidth}{lXX}
\toprule
Tài liệu & Hướng tiếp cận & Nội dung tham khảo \\
\midrule
de Koster và cộng sự \cite{dekoster2007} & Tổng quan hoạt động lấy hàng. & Các quyết định về gom đơn và định tuyến. \\
Gademann--van de Velde \cite{gademann2005} & Gom đơn theo thời gian di chuyển. & Quan hệ giữa nhóm đơn và chi phí lấy hàng. \\
Ratliff--Rosenthal \cite{ratliff1983} & Định tuyến trong kho chữ nhật. & Khai thác cấu trúc lối đi. \\
Briant và cộng sự \cite{briant2023} & Tối ưu tích hợp có hạn hoàn thành. & Mô hình tích hợp và cách xây dựng cận. \\
Shaw \cite{shaw1998} & LNS. & Loại bỏ và tái chèn một phần nghiệm. \\
Ropke--Pisinger \cite{ropke2006} & ALNS. & Lựa chọn toán tử theo hiệu quả quan sát. \\
Mladenović--Hansen \cite{mladenovic1997} & VNS. & Thay đổi cấu trúc lân cận. \\
\bottomrule
\end{tabularx}
\caption{Các hướng giải quyết liên quan và nội dung tham khảo.}\label{tab:related-work}
\end{table}

Trên cơ sở các hướng trong \tabref{tab:related-work}, đề tài cài đặt năm phương pháp: First-Come, First-Served Batching (B0), Greedy Batching (B2), LNS, ALNS và VNS. B0 là mốc kiến tạo đơn giản, gom đơn theo thứ tự đầu vào; B2 sử dụng tiêu chí tham lam để xây dựng nghiệm và làm điểm xuất phát cho ba phương pháp tìm kiếm. So sánh LNS với ALNS nhằm khảo sát lựa chọn toán tử đều và lựa chọn thích ứng trên cùng bộ toán tử. VNS bổ sung một chiến lược đối chứng dựa trên thay đổi cấu trúc lân cận. Cách lựa chọn này cho phép đánh giá từ nghiệm cơ sở đến các cơ chế cải thiện nghiệm, với cùng mô hình và hàm mục tiêu.

Trong cài đặt, B0 dựng tuyến bằng Nearest Neighbor, còn B2 kết hợp Nearest Neighbor với 2-Opt. Vì vậy, so sánh B2 với B0 phản ánh sự khác biệt của cả cách gom đơn và cách định tuyến; không quy toàn bộ chênh lệch chất lượng nghiệm cho riêng bước gom đơn tham lam. Các đối chiếu với B2 cho biết lợi ích của việc tiếp tục tìm kiếm trong ngân sách đã xét, nhưng không bảo đảm một phương pháp luôn tốt hơn trên mọi đầu vào.

\subsection{Mục tiêu và phạm vi đề tài}
Mục tiêu của đề tài là xây dựng và phân tích các thuật toán giải bài toán tích hợp gom đơn, định tuyến và lập lịch, đồng thời đánh giá chất lượng nghiệm và chi phí tính toán của năm phương pháp. Phần thuật toán trình bày nguyên lý, giả mã, ví dụ minh họa và phân tích chi phí của cài đặt. Phần thực nghiệm đối chiếu kết quả trên các nhóm dữ liệu, liên hệ những khác biệt quan sát được với cơ chế tìm kiếm và điều kiện chạy.

Trong phần thực nghiệm, thời gian tính toán là thời gian chương trình chạy để xây dựng và kiểm chứng kế hoạch. Đại lượng này được phân biệt với makespan, là thời gian hoàn tất công việc theo lịch thực hiện trong mô hình kho. Một thuật toán chạy nhanh trên máy tính chưa chắc tạo ra kế hoạch có makespan thấp, và ngược lại.

Các sản phẩm chính gồm bộ giải cài đặt năm phương pháp trên cùng mô hình và bộ đánh giá; cơ chế kiểm chứng tính hợp lệ của nghiệm cùng các chỉ tiêu được xuất ra; và bộ thực nghiệm lưu cấu hình, đầu vào, nghiệm và số liệu để tổng hợp kết quả so sánh. Giao diện mô phỏng hỗ trợ quan sát tuyến đi và lịch thực hiện. Đóng góp của đề tài tập trung vào việc áp dụng, cài đặt và kiểm chứng các phương pháp trong phạm vi bài toán đã xác định, không khẳng định một khung thuật toán mới.

Các câu hỏi đánh giá chính gồm:
\begin{enumerate}
\item Chất lượng nghiệm của B2 khác B0 như thế nào?
\item Tiếp tục tìm kiếm từ nghiệm B2 làm thay đổi chất lượng nghiệm và thời gian tính toán như thế nào?
\item Với cùng bộ toán tử và ngân sách, lựa chọn thích ứng của ALNS cho kết quả khác LNS như thế nào?
\item Chiến lược thay đổi lân cận của VNS cho chất lượng nghiệm, độ ổn định và thời gian tính toán như thế nào so với LNS và ALNS?
\item Quy mô bài toán và ngân sách tìm kiếm ảnh hưởng thế nào đến chất lượng nghiệm và thứ hạng các phương pháp?
\item Thay đổi trọng số mục tiêu tạo ra những đánh đổi nào giữa quãng đường, thời gian hoàn tất và độ trễ?
\end{enumerate}

Phạm vi nghiên cứu là kho tĩnh với số nhân viên cho trước, nhân viên đồng nhất về tốc độ và thời gian thao tác, cùng xe đẩy có sức chứa hữu hạn. Toàn bộ đơn được biết và sẵn sàng tại thời điểm bắt đầu; mỗi đơn được phục vụ trọn vẹn trong một chuyến và hoàn thành khi chuyến quay về depot. Mô hình chưa xét đơn phát sinh trong quá trình vận hành, ca nghỉ hoặc xung đột di chuyển giữa nhân viên.

Thực nghiệm sử dụng các bản đồ kho tổng hợp, nhóm dữ liệu thay đổi quy mô và bộ dữ liệu tham chiếu Kris Small. Kết quả được tổng hợp riêng theo nhóm dữ liệu; các trường hợp nổi bật được lựa chọn để phân tích cùng với bảng kết quả đầy đủ. Việc đánh giá bao gồm kiểm tra tính hợp lệ của nghiệm, chất lượng mục tiêu, các thành phần vận hành, hội tụ và độ ổn định qua các lần chạy. Các kết luận được giới hạn trong mô hình, dữ liệu và ngân sách đã khảo sát; kết quả trên dữ liệu mô phỏng chưa đủ để khẳng định hiệu quả trong mọi kho thực tế hoặc bảo đảm tối ưu toàn cục.
